\chapter{Breking van licht}\label{ch:g10:refraction}

Steek een rietje in een glas water: het lijkt geknakt bij de
waterspiegel. Zwembaden zijn altijd dieper dan ze eruitzien, diamanten
schieten gekleurd vuur uit een kleurloze steen, en de zin die je nu leest
is waarschijnlijk een oceaan overgestoken binnen in een glazen draad
dunner dan een haar. Eén wet verklaart alle vier --- een gelijkheid
tussen twee sinussen, in 1621 gevonden door Snellius, nadat
natuurkundigen vijftien eeuwen lang naar tabellen met hoeken hadden
gestaard zonder het patroon te zien.

\section{Terugkaatsing: een herinnering}

In een homogeen doorzichtig midden plant licht zich rechtlijnig voort ---
de stralen van \cref{ch:g10:light-spectra}. Op de grens tussen twee
middens kaatst een deel van het licht terug (\emph{terugkaatsing}) en
gaat een deel over (\emph{\omterm{def:g10:refraction:refraction}{breking}}). Beide gehoorzamen aan exacte
meetkundige wetten, met één afspraak: hoeken worden gemeten vanaf de
\emph{\omterm{def:g10:refraction:angles}{normaal}}, nooit vanaf het oppervlak.

\begin{definition}[Normaal en invalshoek]\label{def:g10:refraction:angles}
Op het punt waar een straal een oppervlak raakt, is de
\emph{normaal}\index{normaal} de lijn die daar loodrecht op het oppervlak
staat. De \emph{invalshoek}\index{invalshoek} $i_1$ is de hoek tussen de
invallende straal en de normaal; de terugkaatsingshoek en de brekingshoek
worden op dezelfde manier vanaf de normaal gemeten.
\end{definition}

\begin{proposition}[Wet van de terugkaatsing]\label{prop:g10:refraction:reflection}
De teruggekaatste straal ligt in het vlak door de invallende straal en de
\omterm{def:g10:refraction:angles}{normaal}, aan de andere kant van de \omterm{def:g10:refraction:angles}{normaal}, en de terugkaatsingshoek is
gelijk aan de \omterm{def:g10:refraction:angles}{invalshoek}:
\[
i_1' = i_1 .
\]
\end{proposition}
\admitted

\begin{example}[Hoeken juist aflezen]\label{ex:g10:refraction:mirror}
Een laserbundel treft een spiegel ``onder $35^\circ$'' --- gemeten vanaf
het spiegeloppervlak, zoals bundels vaak worden beschreven. De \omterm{def:g10:refraction:angles}{invalshoek}
is dan $90^\circ - 35^\circ = 55^\circ$, dus de teruggekaatste straal
vertrekt ook onder $55^\circ$ met de \omterm{def:g10:refraction:angles}{normaal}, en de bundel draait over
$180^\circ - 2 \times 55^\circ = 70^\circ$. De afspraak over de \omterm{def:g10:refraction:angles}{normaal}
vergeten is de meest gemaakte fout van dit hoofdstuk; ze kost elke keer
$90^\circ$.
\end{example}

\section{Breking en de brekingsindex}

Houd een potlood in een halfvol glas: het lijkt gebroken aan het
oppervlak. Licht uit het deel onder water verandert van richting bij het
verlaten van het water, en dus reconstrueert het oog --- dat rechte
stralen veronderstelt --- de punt op een plaats waar die niet is.

\begin{definition}[Breking]\label{def:g10:refraction:refraction}
\emph{Breking}\index{breking} is de richtingsverandering van licht bij het
oversteken van de grens tussen twee doorzichtige middens. De straal in
het tweede midden is de \emph{gebroken straal}\index{gebroken straal}, en
haar hoek met de \omterm{def:g10:refraction:angles}{normaal} is de brekingshoek $i_2$.
\end{definition}

\omterm{def:g10:refraction:refraction}{Breking} treedt op doordat licht in verschillende middens verschillende
snelheden heeft; elk materiaal wordt gekenmerkt door hoeveel het het
licht afremt.

\begin{definition}[Brekingsindex]\label{def:g10:refraction:index}
De \emph{brekingsindex}\index{brekingsindex} van een doorzichtig midden
is de verhouding
\[
n = \frac{c}{v},
\]
waarin $c = \qty{3.00e8}{m/s}$ de lichtsnelheid in vacuüm is en $v$ de
lichtsnelheid in het midden. Omdat $v \leq c$, geldt $n \geq 1$; de index
heeft geen \omterm{def:g10:orders-of-magnitude:unit}{eenheid}. Enkele bruikbare waarden:
\begin{center}
\begin{tabular}{lccccc}
midden & lucht & water & plexiglas & glas & diamant \\
$n$ & $1.0003$ & $1.33$ & $1.49$ & $1.50$ & $2.42$ \\
\end{tabular}
\end{center}
\end{definition}

\begin{example}[Lichtsnelheid in water]\label{ex:g10:refraction:speed}
In water is $v = \dfrac{c}{n} = \dfrac{\qty{3.00e8}{m/s}}{1.33}
= \qty{2.26e8}{m/s}$. In diamant kruipt het licht met
$\qty{3.00e8}{m/s}/2.42 = \qty{1.24e8}{m/s}$ --- minder dan de helft van
zijn vacuümsnelheid. Voor lucht verandert $n = 1.0003$ de snelheid met
$0.03\%$: in dit hoofdstuk behandelen we lucht als vacuüm en nemen we
$n_{\text{lucht}} = 1.00$.
\end{example}

\begin{theorem}[Brekingswet van Snellius]\label{thm:g10:refraction:snell}
\index{wet van Snellius}
De \omterm{def:g10:refraction:refraction}{gebroken straal} ligt in het vlak van de invallende straal en de
\omterm{def:g10:refraction:angles}{normaal}, aan de andere kant van de \omterm{def:g10:refraction:angles}{normaal}, en de \omterm{def:g10:refraction:angles}{invalshoek} en de
brekingshoek voldoen aan
\[
n_1 \sin i_1 = n_2 \sin i_2 ,
\]
waarin $n_1$ en $n_2$ de brekingsindices van de twee middens zijn. De weg
is omkeerbaar: licht dat de \omterm{def:g10:refraction:refraction}{gebroken straal} terugloopt, verlaat het
oppervlak langs de invallende straal.
\end{theorem}
\admitted

\begin{remark}\label{rem:g10:refraction:honest}
De wetten van terugkaatsing en \omterm{def:g10:refraction:refraction}{breking} staan hier als experimentele
feiten. Ze worden eerlijk afgeleid in het volume van bachelorjaar 1,
waar licht als golf wordt behandeld: beide wetten --- en de formule
$n = c/v$ zelf --- volgen dan uit het afremmen van de golf in het
dichtere midden.
\end{remark}

\begin{omfigure}
\begin{tikzpicture}[scale=1.0, font=\small]
  \fill[blue!8] (-3.2,-2.1) rectangle (3.2,0);
  \draw[thick] (-3.2,0) -- (3.2,0);
  \draw[dashed] (0,-2.1) -- (0,2.3);
  \coordinate (O) at (0,0);
  \coordinate (I) at (-1.8,2.0);
  \coordinate (R) at (1.8,2.0);
  \coordinate (T) at (1.0,-1.73);
  \coordinate (N) at (0,2.3);
  \coordinate (Nd) at (0,-2.1);
  \draw[-Stealth, omDef, thick] (I) -- (O);
  \draw[-Stealth, omExo, thick] (O) -- (R);
  \draw[-Stealth, omThm, thick] (O) -- (T);
  \pic[draw, angle radius=9mm, "$i_1$", angle eccentricity=1.35] {angle = N--O--I};
  \pic[draw, angle radius=7mm, "$i_1'$", angle eccentricity=1.4] {angle = R--O--N};
  \pic[draw, angle radius=9mm, "$i_2$", angle eccentricity=1.35] {angle = Nd--O--T};
  \node[anchor=west] at (-3.1,0.45) {midden 1 ($n_1$)};
  \node[anchor=west] at (-3.1,-0.55) {midden 2 ($n_2$)};
\end{tikzpicture}

{\small Invallende (blauw), teruggekaatste (grijs) en gebroken (rood)
straal. Alle hoeken worden vanaf de \omterm{def:g10:refraction:angles}{normaal} (gestippeld) gemeten; bij het
binnengaan van het dichtere midden ($n_2 > n_1$) buigt de straal
\emph{naar} de \omterm{def:g10:refraction:angles}{normaal} toe.}
\end{omfigure}

\begin{example}[Van lucht naar water]\label{ex:g10:refraction:airwater}
Een straal treft een rustige vijver onder $i_1 = 30^\circ$. Met
$n_1 = 1.00$ en $n_2 = 1.33$:
\[
\sin i_2 = \frac{n_1 \sin i_1}{n_2} = \frac{\sin 30^\circ}{1.33}
= \frac{0.500}{1.33} = 0.376,
\qquad
i_2 = 22.1^\circ .
\]
De straal buigt naar de \omterm{def:g10:refraction:angles}{normaal} toe, zoals altijd bij het binnengaan van
een trager midden ($n_2 > n_1$). Bij het verlaten van het water loopt
dezelfde berekening achterstevoren en buigt de straal juist \emph{van}
de \omterm{def:g10:refraction:angles}{normaal} \emph{af}.
\end{example}

\begin{method}[Een brekingsopgave oplossen]\label{met:g10:refraction:method}
\begin{enumerate}
  \item Teken het oppervlak, de \omterm{def:g10:refraction:angles}{normaal} en de straal; bepaal $n_1$ (het
        midden van de invallende straal) en $n_2$.
  \item Ga na dat de gegeven hoeken vanaf de \omterm{def:g10:refraction:angles}{normaal} gemeten zijn; zet ze
        om als ze vanaf het oppervlak gemeten zijn.
  \item Schrijf $n_1 \sin i_1 = n_2 \sin i_2$ en zonder de onbekende
        sinus af; haal de hoek eruit met de inverse sinus op de
        rekenmachine, zoals in het wiskundevolume.
  \item Controleer: de straal buigt naar de \omterm{def:g10:refraction:angles}{normaal} toe als $n_2 > n_1$,
        ervan af als $n_2 < n_1$, en $\sin i_2$ moet tussen $0$ en $1$
        uitkomen.
\end{enumerate}
\end{method}

Het experiment achter de wet past op een tafel: een halve schijf
plexiglas, een gradenboog, een laser. Ruwe hoekparen $(i_1, i_2)$ ogen
grillig --- Ptolemaeus zette ze rond het jaar 150 in een tabel en
concludeerde ten onrechte dat $i_2$ evenredig was met $i_1$. Zet in
plaats daarvan sinus tegen sinus uit, en de meetpunten strekken zich tot
een rechte door de oorsprong met richtingscoëfficiënt $n_1/n_2$: die
rechte \emph{is} de wet van Snellius.

\begin{omfigure}
\begin{tikzpicture}
\begin{axis}[omaxis, width=7.2cm, height=5.4cm,
  xlabel={$\sin i_1$}, ylabel={$\sin i_2$},
  xmin=0, xmax=1.08, ymin=0, ymax=0.88,
  xtick={0,0.2,0.4,0.6,0.8,1.0}, ytick={0.2,0.4,0.6,0.8}]
  \addplot[omThm, thick, domain=0:1.02] {x/1.33};
  \addplot[only marks, mark=*, mark size=1.6pt, omDef] coordinates
    {(0.174,0.131) (0.342,0.257) (0.500,0.376) (0.643,0.483)
     (0.766,0.576) (0.866,0.651) (0.940,0.707)};
\end{axis}
\end{tikzpicture}

{\small Metingen van lucht naar water voor $i_1 = 10^\circ$ tot
$70^\circ$: sinus tegen sinus uitgezet liggen de punten op een rechte
door de oorsprong met richtingscoëfficiënt $1/1.33$.}
\end{omfigure}

\begin{proposition}[Schijnbare diepte]\label{prop:g10:refraction:depth}
\index{schijnbare diepte}
Een voorwerp op diepte $d$ onder water, van boven bekeken onder kleine
hoeken, lijkt op de \emph{\omterm{prop:g10:refraction:depth}{schijnbare diepte}}
\[
d' = \frac{d}{n}
\]
te liggen, waarin $n$ de index van water is. Een bad van \qty{1.20}{m}
diep lijkt maar $1.20/1.33 = \qty{0.90}{m}$ diep --- een klassieke
oorzaak van onvoorzichtige duiksprongen.
\end{proposition}
\admitted

\begin{remark}\label{rem:g10:refraction:depthhonest}
De formule volgt uit de wet van Snellius toegepast op de bijna verticale
stralen die het oog bereiken; die afleiding staat in het volume van
bachelorjaar 1. De richting van het effect heeft geen formule nodig:
stralen die het water verlaten buigen van de \omterm{def:g10:refraction:angles}{normaal} af, dus verlengt het
oog ze terug naar een punt hoger dan het voorwerp.
\end{remark}

\section{Dispersie: één wet per kleur}

Een prisma maakt van wit zonlicht een regenboog, zoals de spectra van
\cref{ch:g10:light-spectra} lieten zien. De \omterm{def:g10:refraction:refraction}{breking} verklaart
\emph{waarom}: de \omterm{def:g10:refraction:index}{brekingsindex} van een midden is niet helemaal
dezelfde voor elke kleur.

\begin{definition}[Dispersie]\label{def:g10:refraction:dispersion}
Een midden is \emph{dispersief}\index{dispersie} wanneer zijn
\omterm{def:g10:refraction:index}{brekingsindex} van de \omterm{def:g10:light-spectra:wavelength}{golflengte} van het licht afhangt. In glas en water
is $n$ iets groter voor blauw licht (korte \omterm{def:g10:light-spectra:wavelength}{golflengte}) dan voor rood
licht (lange \omterm{def:g10:light-spectra:wavelength}{golflengte}) --- en dus buigt blauw sterker af.
\end{definition}

\begin{example}[Twee kleuren, één oppervlak]\label{ex:g10:refraction:twocolors}
Een bundel \omterm{def:g10:light-spectra:white}{wit licht} treft kroonglas onder $i_1 = 60^\circ$. Voor dit
glas is $n_{\text{rood}} = 1.51$ en $n_{\text{blauw}} = 1.53$. De wet van
Snellius, één keer per kleur:
\[
\sin i_{\text{rood}} = \frac{\sin 60^\circ}{1.51} = 0.574,
\quad i_{\text{rood}} = 35.0^\circ ;
\qquad
\sin i_{\text{blauw}} = \frac{\sin 60^\circ}{1.53} = 0.566,
\quad i_{\text{blauw}} = 34.5^\circ .
\]
Eén oppervlak splitst \omterm{def:g10:light-spectra:white}{wit licht} over een halve graad --- onzichtbaar op
een ruit, maar een prisma breekt twee keer in dezelfde draaizin,
verbreedt de waaier, en een paar \omterm{def:g10:orders-of-magnitude:unit}{meter} verderop liggen de kleuren
centimeters uit elkaar.
\end{example}

\begin{remark}\label{rem:g10:refraction:rainbow}
De regenboog is dezelfde berekening, uitgevoerd binnen in een
regendruppel: zonlicht breekt bij het binnengaan van de druppel, kaatst
terug tegen de achterwand en breekt bij het naar buiten gaan opnieuw. De
index van water loopt van $1.331$ (rood) tot $1.343$ (blauw), dus verlaat
elke kleur de druppel onder haar eigen hoek --- en geen twee waarnemers
zien dezelfde regenboog.
\end{remark}

\section{Totale terugkaatsing}

De wet van Snellius bevat een valstrik. Van water naar lucht buigt de
\omterm{def:g10:refraction:refraction}{gebroken straal} van de \omterm{def:g10:refraction:angles}{normaal} af: $\sin i_2 = 1.33 \sin i_1 > \sin i_1$.
Vergroot $i_1$, en de formule eist al snel $\sin i_2 > 1$, wat geen
enkele hoek kan leveren. Dan houdt de \omterm{def:g10:refraction:refraction}{breking} eenvoudigweg op te bestaan.

\begin{definition}[Grenshoek en totale terugkaatsing]\label{def:g10:refraction:critical}
Licht plant zich voort in een midden met index $n_1$ en valt op een
midden met kleinere index $n_2 < n_1$. De \emph{grenshoek}\index{grenshoek}
$i_c$ is de \omterm{def:g10:refraction:angles}{invalshoek} waarvoor de \omterm{def:g10:refraction:refraction}{gebroken straal} langs het oppervlak
scheert ($i_2 = 90^\circ$). Voor $i_1 > i_c$ bestaat er geen \omterm{def:g10:refraction:refraction}{gebroken
straal}: het oppervlak kaatst al het licht terug het eerste midden in,
volgens de wet van de terugkaatsing. Dit is \emph{totale
terugkaatsing}\index{totale terugkaatsing}.
\end{definition}

\begin{proposition}[Formule voor de grenshoek]\label{prop:g10:refraction:criticalformula}
Voor $n_1 > n_2$ geldt
\[
\sin i_c = \frac{n_2}{n_1} .
\]
\end{proposition}

\begin{proof}
Stel $i_2 = 90^\circ$ in de wet van Snellius: $n_1 \sin i_c =
n_2 \sin 90^\circ = n_2$. Voor $i_1 > i_c$ zou de wet van Snellius
$\sin i_2 = \frac{n_1}{n_2}\sin i_1 > 1$ vergen: er bestaat geen gebroken
richting. (Dat het licht dan \emph{volledig} wordt teruggekaatst --- en
niet geabsorbeerd --- is een experimenteel feit, bevestigd door de
golftheorie uit het volume van bachelorjaar 1.)
\end{proof}

\begin{omfigure}
\begin{tikzpicture}[scale=0.92, font=\small]
  \fill[blue!8] (-3.6,-1.9) rectangle (3.6,0);
  \draw[thick] (-3.6,0) -- (3.6,0);
  \draw[dashed] (-2.2,-1.6) -- (-2.2,1.6);
  \draw[dashed] (0.3,-1.6) -- (0.3,1.6);
  \draw[dashed] (2.2,-1.6) -- (2.2,1.6);
  \draw[-Stealth, omDef, thick] (-2.95,-1.30) -- (-2.2,0);
  \draw[-Stealth, omThm, thick] (-2.2,0) -- (-1.20,1.12);
  \draw[-Stealth, omDef, thick] (-0.83,-0.99) -- (0.3,0);
  \draw[-Stealth, omThm, thick] (0.3,0) -- (1.75,0.06);
  \draw[-Stealth, omDef, thick] (0.90,-0.75) -- (2.2,0);
  \draw[-Stealth, omDef, thick] (2.2,0) -- (3.50,-0.75);
  \node at (-2.2,-1.85) {$i_1 < i_c$};
  \node at (0.3,-1.85) {$i_1 = i_c$};
  \node at (2.2,-1.85) {$i_1 > i_c$};
\end{tikzpicture}

{\small Van water naar lucht, bij toenemende inval: \omterm{def:g10:refraction:refraction}{breking} van de
\omterm{def:g10:refraction:angles}{normaal} af, dan een \omterm{def:g10:refraction:refraction}{gebroken straal} die bij de \omterm{def:g10:refraction:critical}{grenshoek} langs het
oppervlak scheert, dan \omterm{def:g10:refraction:critical}{totale terugkaatsing} --- het oppervlak is een
volmaakte spiegel geworden.}
\end{omfigure}

\begin{example}[Drie grenshoeken]\label{ex:g10:refraction:criticalvalues}
Naar lucht toe: voor water is $\sin i_c = 1/1.33 = 0.752$, dus
$i_c = 48.8^\circ$; voor glas is $\sin i_c = 1/1.50$, dus
$i_c = 41.8^\circ$; voor diamant is $\sin i_c = 1/2.42 = 0.413$, dus
$i_c = 24.4^\circ$. Hoe hoger de index, hoe smaller de ontsnappingskegel:
in diamant zit elke straal die meer dan $24.4^\circ$ van de \omterm{def:g10:refraction:angles}{normaal} van
een facet afwijkt gevangen --- het geheim van de schittering, ontleed in
\cref{pb:g10:refraction:1}.
\end{example}

\section{Glasvezels}

\omterm{def:g10:refraction:critical}{Totale terugkaatsing} maakt van een glazen draad een buis voor licht: een
straal die bijna evenwijdig aan de as binnenkomt, treft de wand ver
voorbij de \omterm{def:g10:refraction:critical}{grenshoek}, kaatst zonder verlies terug, en doet dat miljoenen
keren per kilometer zonder te lekken.

\begin{definition}[Glasvezel]\label{def:g10:refraction:fiber}
Een \emph{glasvezel}\index{glasvezel} is een dunne glazen draad die
bestaat uit een \emph{kern}\index{kern} met index $n_1$, gehuld in een
\emph{mantel}\index{mantel} met een iets kleinere index $n_2 < n_1$.
Licht dat aan het ene uiteinde wordt ingebracht, reist door de kern door
herhaalde \omterm{def:g10:refraction:critical}{totale terugkaatsing} op de grens tussen kern en mantel.
\end{definition}

\begin{omfigure}
\begin{tikzpicture}[scale=1.0, font=\small]
  \draw[omExo, thick] (0,0.85) -- (7,0.85);
  \draw[omExo, thick] (0,-0.85) -- (7,-0.85);
  \fill[blue!8] (0,-0.5) rectangle (7,0.5);
  \draw[thick] (0,0.5) -- (7,0.5);
  \draw[thick] (0,-0.5) -- (7,-0.5);
  \draw[-Stealth, omThm, thick] (0,-0.36) -- (1.0,0.5) -- (2.16,-0.5)
    -- (3.32,0.5) -- (4.48,-0.5) -- (5.64,0.5) -- (6.6,-0.33);
  \node[anchor=south] at (3.5,0.9) {mantel ($n_2 = 1.46$)};
  \node[fill=blue!8, inner sep=2pt] at (3.5,-0.25) {kern ($n_1 = 1.48$)};
\end{tikzpicture}

{\small Licht geleid langs de \omterm{def:g10:refraction:fiber}{kern} van een vezel door \omterm{def:g10:refraction:critical}{totale
terugkaatsing} (de kaatshoeken zijn sterk overdreven: echte stralen
scheren onder meer dan $80^\circ$ met de \omterm{def:g10:refraction:angles}{normaal} langs de wand).}
\end{omfigure}

\begin{example}[Een vezel in getallen]\label{ex:g10:refraction:fibernumbers}
Met $n_1 = 1.48$ en $n_2 = 1.46$:
\[
\sin i_c = \frac{1.46}{1.48} = 0.986,
\qquad i_c = 80.6^\circ :
\]
alleen stralen die minder dan ongeveer $9^\circ$ van de as afwijken
worden geleid --- wat prima is, want precies zo wordt laserlicht
ingebracht. Binnen de \omterm{def:g10:refraction:fiber}{kern} reist het licht met
$v = c/1.48 = \qty{2.03e8}{m/s}$: een puls doorkruist een vezel van
\qty{100}{km} in ongeveer een halve milliseconde.
\end{example}

\begin{remark}[Ordegrootte van dataverbindingen]\label{rem:g10:refraction:datarates}
Zeekabels --- bundels van een paar dozijn vezels op de oceaanbodem ---
dragen vrijwel al het intercontinentale verkeer. Eén moderne vezel
vervoert in de orde van $10^{13}$ bits per seconde door vele \omterm{def:g10:light-spectra:wavelength}{golflengten}
tegelijk te versturen; een hele kabel haalt $10^{14}$ bits per seconde,
met om de \qty{80}{km} een versterker. Satellieten verwerken ter
vergelijking een fractie van een procent van het verkeer.
\end{remark}

\section{Oefeningen}

\begin{exercise}[$\star$]\label{exo:g10:refraction:1}
Een laserbundel treft een vlakke spiegel onder $25^\circ$ \emph{met het
spiegeloppervlak}. Geef de \omterm{def:g10:refraction:angles}{invalshoek}, de terugkaatsingshoek en de hoek
tussen de invallende en de teruggekaatste bundel.
\end{exercise}

\begin{exercise}[$\star$]\label{exo:g10:refraction:2}
\begin{enumerate}
  \item Bereken de lichtsnelheid in water ($n = 1.33$) en in diamant
        ($n = 2.42$).
  \item In een bepaalde vaste stof reist licht met \qty{1.97e8}{m/s}.
        Bereken de \omterm{def:g10:refraction:index}{brekingsindex} en benoem het materiaal.
\end{enumerate}
\end{exercise}

\begin{exercise}[$\star$]\label{exo:g10:refraction:3}
Een straal in lucht treft een glazen blok ($n = 1.50$) onder
$i_1 = 40^\circ$. Bereken de brekingshoek. Onder welke hoek zou een
straal binnen in het glas het oppervlak moeten treffen om onder
$40^\circ$ de lucht in te gaan?
\end{exercise}

\begin{exercise}[$\star$]\label{exo:g10:refraction:4}
Zonlicht bereikt het oppervlak van een meer onder $55^\circ$ met de
verticaal. Bereken de hoek van de \omterm{def:g10:refraction:refraction}{gebroken straal} en zeg of de straal
naar de \omterm{def:g10:refraction:angles}{normaal} toe of ervan af buigt --- met de reden in één woord.
\end{exercise}

\begin{exercise}[$\star$]\label{exo:g10:refraction:5}
Om een heldere vloeistof te herkennen wordt een straal onder
$i_1 = 45^\circ$ op haar oppervlak gestuurd; de \omterm{def:g10:refraction:refraction}{gebroken straal} wordt
gemeten op $i_2 = 32^\circ$. Bereken de \omterm{def:g10:refraction:index}{brekingsindex} en herken de
vloeistof aan de tabel van \cref{def:g10:refraction:index}.
\end{exercise}

\begin{exercise}[$\star\star$]\label{exo:g10:refraction:6}
Een halve schijf van een doorzichtige kunststof levert de volgende
metingen:
\begin{center}
\begin{tabular}{lccccc}
$i_1$ & $20^\circ$ & $30^\circ$ & $40^\circ$ & $50^\circ$ & $60^\circ$ \\
$i_2$ & $13.5^\circ$ & $19.5^\circ$ & $25.5^\circ$ & $31.0^\circ$
      & $35.5^\circ$ \\
\end{tabular}
\end{center}
\begin{enumerate}
  \item Bereken $\dfrac{\sin i_1}{\sin i_2}$ voor elk paar. Wat valt je
        op?
  \item Leid de \omterm{def:g10:refraction:index}{brekingsindex} van de kunststof af en benoem haar.
  \item Waarom is $\sin i_2$ tegen $\sin i_1$ uitzetten een betere toets
        van de wet van Snellius dan $i_2$ tegen $i_1$ uitzetten?
\end{enumerate}
\end{exercise}

\begin{exercise}[$\star\star$]\label{exo:g10:refraction:7}
Bereken de \omterm{def:g10:refraction:critical}{grenshoek} naar lucht toe voor water, glas ($n = 1.50$) en
diamant. Leg daarna met de wet van Snellius uit waarom er geen \omterm{def:g10:refraction:critical}{grenshoek}
bestaat voor licht dat \emph{vanuit} lucht \emph{in} water gaat.
\end{exercise}

\begin{exercise}[$\star\star$]\label{exo:g10:refraction:8}
\begin{enumerate}
  \item Een zwembad is \qty{2.0}{m} diep. Welke diepte neemt een
        zwemmer aan de rand waar?
  \item Een reiger loert op een vis die \qty{40}{cm} onder het oppervlak
        zwemt. Op welke diepte lijkt de vis te zitten, en moet de reiger
        boven, op of onder het beeld toeslaan?
\end{enumerate}
\end{exercise}

\begin{exercise}[$\star\star$]\label{exo:g10:refraction:9}
Een straal treft een ruit ($n = 1.50$, evenwijdige zijden) onder
$50^\circ$.
\begin{enumerate}
  \item Bereken de brekingshoek aan het eerste vlak.
  \item De \omterm{def:g10:refraction:refraction}{gebroken straal} treft het tweede vlak van binnenuit, onder
        dezelfde hoek. Bereken de uittredehoek en vergelijk die met
        $50^\circ$.
  \item Wat doet een dikke ruit dus met de richting en met de plaats van
        een straal die haar doorkruist?
\end{enumerate}
\end{exercise}

\begin{exercise}[$\star\star$]\label{exo:g10:refraction:10}
\omterm{def:g10:light-spectra:white}{Wit licht} treft een blok flintglas onder $55^\circ$. Voor dit glas is
$n_{\text{rood}} = 1.61$ en $n_{\text{blauw}} = 1.66$. Bereken de
brekingshoeken van de rode en de blauwe straal en de hoek daartussen.
Welke kleur komt het dichtst bij de \omterm{def:g10:refraction:angles}{normaal} terecht?
\end{exercise}

\begin{exercise}[$\star\star$]\label{exo:g10:refraction:11}
Een \omterm{def:g10:refraction:fiber}{glasvezel} heeft een \omterm{def:g10:refraction:fiber}{kern} met index $1.48$ en een \omterm{def:g10:refraction:fiber}{mantel} met index
$1.46$.
\begin{enumerate}
  \item Bereken de \omterm{def:g10:refraction:critical}{grenshoek} op de grens tussen \omterm{def:g10:refraction:fiber}{kern} en \omterm{def:g10:refraction:fiber}{mantel}.
  \item Een geleide straal kaatst precies onder de \omterm{def:g10:refraction:critical}{grenshoek}. Toon aan
        dat hij over \qty{1.0}{km} rechte vezel ongeveer \qty{14}{m}
        meer aflegt dan een straal die recht langs de as loopt.
\end{enumerate}
\end{exercise}

\begin{exercise}[$\star\star\star$]\label{exo:g10:refraction:12}
Het oog van een duiker zit \qty{2.0}{m} onder een volmaakt rustig
oppervlak. Door de \omterm{def:g10:refraction:refraction}{breking} ziet de duiker de \emph{hele} hemel
samengeperst tot een heldere schijf recht boven zich --- het venster van
Snellius.
\begin{enumerate}
  \item Leg uit waarom licht van de horizon het oog van de duiker onder
        de \omterm{def:g10:refraction:critical}{grenshoek} van water, $48.8^\circ$ met de verticaal, bereikt.
  \item Bereken de straal van de heldere schijf aan het oppervlak, met de
        tangens in de rechthoekige driehoek oog--verticaal--schijfrand.
  \item Wat ziet de duiker aan het oppervlak \emph{buiten} de schijf?
\end{enumerate}
\end{exercise}

\begin{exercise}[$\star\star\star$]\label{exo:g10:refraction:13}
Een prisma heeft hoeken van $30^\circ$, $60^\circ$ en $90^\circ$. Een
dunne bundel \omterm{def:g10:light-spectra:white}{wit licht} komt loodrecht binnen door een van de zijden naast
de rechte hoek, gaat onafgebogen door en treft het lange vlak van
binnenuit onder $i_1 = 30^\circ$. Voor dit glas is $n_{\text{rood}} =
1.51$ en $n_{\text{blauw}} = 1.53$.
\begin{enumerate}
  \item Bereken de uittredehoeken van de rode en de blauwe straal en de
        hoekbreedte van de waaier daartussen.
  \item Het prisma wordt vervangen door een prisma met hoeken
        $45^\circ$--$45^\circ$--$90^\circ$, zodat de inwendige \omterm{def:g10:refraction:angles}{invalshoek}
        $45^\circ$ wordt. Toon aan dat er helemaal geen licht door het
        lange vlak naar buiten komt.
\end{enumerate}
\end{exercise}

\begin{exercise}[$\star\star\star$]\label{exo:g10:refraction:14}
Periscopen vouwen licht met glazen prisma's van
$45^\circ$--$45^\circ$--$90^\circ$ ($n = 1.50$): de bundel komt loodrecht
door een korte zijde binnen en treft het lange vlak van binnenuit onder
$45^\circ$.
\begin{enumerate}
  \item Toon aan dat het lange vlak als een volmaakte spiegel werkt.
        Waarom is dat beter dan een metalen spiegel, die bij elke
        terugkaatsing een paar procent absorbeert?
  \item De periscoop loopt vol en het lange vlak raakt nu water.
        Bereken de nieuwe \omterm{def:g10:refraction:critical}{grenshoek} van de grens glas--water en toon aan
        dat de spiegel het begeeft.
  \item Onder welke hoek verlaat het licht van vraag 2 het prisma en gaat
        het het water in?
\end{enumerate}
\end{exercise}

\begin{exercise}[$\star\star\star$]\label{exo:g10:refraction:15}
Een transatlantische kabel loopt over \qty{6200}{km} vezel (index van de
\omterm{def:g10:refraction:fiber}{kern} $1.47$) tussen Londen en New York.
\begin{enumerate}
  \item Bereken de lichtsnelheid in de \omterm{def:g10:refraction:fiber}{kern} en de reistijd van een puls
        in één richting.
  \item Een radiogolf reist met $c$. Hoe lang zou die over dezelfde
        afstand doen, en hoeveel heen-en-weertijd kost het glas van de
        vezel vergeleken met radio?
  \item Sommige financiële bedrijven betalen fortuinen om milliseconden
        van deze verbinding af te snoepen. Stel twee verschillende
        natuurkundige manieren voor om het licht eerder te laten
        aankomen.
\end{enumerate}
\end{exercise}

\section{Probleem: De diamant, het zwembad en de vezel}

\begin{problem}[{Weekendopgave --- drie levens van de wet van Snellius:
een zwembad dat over zijn diepte liegt, een steen die zo geslepen is dat
licht er niet uit kan, en de oceaan van zestig milliseconden onder het
internet}]
\label{pb:g10:refraction:1}
Eén gelijkheid tussen twee sinussen, drie loopbanen. In een zwembad
bedriegt ze je ogen over de diepte en perst ze de hele hemel in een
cirkel. In een diamant, tot het uiterste gedreven, weigert ze het licht
naar buiten te laten --- en juweliers slijpen hun stenen al eeuwen rond
die weigering. Onder de Atlantische Oceaan leidt ze laserpulsen binnen in
een sliert glas, en legt ze de maximumsnelheid van het wereldwijde
internet vast. Deze opgave werkt de drie levens op volgorde af, met
telkens dezelfde wet.

\medskip
\noindent\textbf{Deel I --- Opwarmen.}
\begin{enumerate}
  \item Bereken de lichtsnelheid in water ($n = 1.33$) en in diamant
        ($n = 2.42$). Met welke factor remt diamant het licht af?
  \item Een straal in lucht treft een vijver onder $i_1 = 40^\circ$.
        Bereken de brekingshoek.
  \item Een straal in het water treft het oppervlak van onderaf onder
        $28.9^\circ$. Onder welke hoek gaat hij de lucht in? Noem het
        principe dat dit illustreert.
  \item Bereken de \omterm{def:g10:refraction:critical}{grenshoek} van de grens water--lucht, en beschrijf wat
        er gebeurt met een straal onder water die het oppervlak onder
        $60^\circ$ treft.
  \item De index van lucht is $1.0003$. Bereken de lichtsnelheid in
        lucht en verantwoord de benadering $n_{\text{lucht}} = 1.00$ die
        overal wordt gebruikt.
\end{enumerate}

\medskip
\noindent\textbf{Deel II --- Het zwembad.}
\begin{enumerate}[resume]
  \item Een zwembad is \qty{3.0}{m} diep. Bereken met de formule voor de
        \omterm{prop:g10:refraction:depth}{schijnbare diepte} (\cref{prop:g10:refraction:depth}) welke diepte
        iemand die erboven staat waarneemt.
  \item Leg met een schets van twee stralen (in woorden beschreven) uit
        waarom een muntstuk op de bodem omhoog lijkt te komen: wat doen
        stralen die het water verlaten, en waar plaatst het oog hun
        snijpunt?
  \item Het oog van een duiker zit op \qty{3.0}{m}. Bereken de straal van
        het venster van Snellius --- de heldere schijf waardoor de duiker
        de hele hemel ziet (\cref{exo:g10:refraction:12}).
  \item Een lamp op de bodem stuurt stralen omhoog onder $30^\circ$,
        $45^\circ$ en $50^\circ$ met de verticaal. Zeg voor elke straal
        of hij ontsnapt en onder welke hoek, of wat er anders gebeurt.
  \item Een harpoenvisser op de oever mikt op een vis. Moet hij boven,
        op of onder het beeld van de vis mikken --- en ziet de vis, die
        terugkijkt, de visser ook verplaatst?
\end{enumerate}

\medskip
\noindent\textbf{Deel III --- De diamant.}
\begin{enumerate}[resume]
  \item Bereken de \omterm{def:g10:refraction:critical}{grenshoek} van diamant ($n = 2.42$) naar lucht toe, en
        vergelijk die met die van glas ($n = 1.52$), namelijk
        $41.1^\circ$.
  \item Een straal binnen in de steen treft een facet onder $30^\circ$
        met de \omterm{def:g10:refraction:angles}{normaal} ervan. Wat gebeurt er in de diamant? En in de
        glazen imitatie? Leg in één zin uit waarom de slijpvorm van een
        diamant, samen met haar kleine \omterm{def:g10:refraction:critical}{grenshoek}, de schittering
        oplevert.
  \item Diamant is ook sterk \omterm{def:g10:refraction:dispersion}{dispersief}: $n_{\text{rood}} = 2.41$,
        $n_{\text{blauw}} = 2.45$. \omterm{def:g10:light-spectra:white}{Wit licht} komt onder $45^\circ$ een
        facet binnen; bereken de brekingshoek voor elke kleur en de
        hoeksplitsing.
  \item Verderop bereiken de rode en de blauwe straal allebei van
        binnenuit een uittredefacet onder $17.0^\circ$. Bereken de twee
        uittredehoeken en de splitsing van de uittredende waaier.
        Vergelijk met de splitsing bij binnenkomst: wat gebeurt er met
        de \omterm{def:g10:refraction:dispersion}{dispersie} in de buurt van de \omterm{def:g10:refraction:critical}{grenshoek}?
  \item Een druppel water ($n = 1.33$) bedekt een facet. Bereken de
        nieuwe \omterm{def:g10:refraction:critical}{grenshoek} van de grens diamant--water, bepaal het lot van
        de straal van $30^\circ$ uit vraag 12, en leg uit waarom
        juweliers hun diamanten angstvallig schoon houden.
\end{enumerate}

\medskip
\noindent\textbf{Deel IV --- De vezel.}
\begin{enumerate}[resume]
  \item Een vezel heeft een \omterm{def:g10:refraction:fiber}{kern} met index $1.48$ en een \omterm{def:g10:refraction:fiber}{mantel} met index
        $1.46$. Bereken de \omterm{def:g10:refraction:critical}{grenshoek} op de grens tussen \omterm{def:g10:refraction:fiber}{kern} en \omterm{def:g10:refraction:fiber}{mantel}.
  \item Een kale glasdraad in lucht zou een \omterm{def:g10:refraction:critical}{grenshoek} van $42.5^\circ$
        hebben --- ogenschijnlijk veel beter. Geef twee redenen waarom
        de \omterm{def:g10:refraction:fiber}{mantel} er toch is. (Denk aan stof, krassen en aan wat aanraken
        van het oppervlak met de \omterm{def:g10:refraction:critical}{totale terugkaatsing} zou doen.)
  \item Bereken de lichtsnelheid in de \omterm{def:g10:refraction:fiber}{kern}, en daarna de reistijd van
        een puls in één richting over \qty{6000}{km} vezel.
  \item Een straal die precies onder de \omterm{def:g10:refraction:critical}{grenshoek} kaatst, legt
        $1/\sin i_c$ keer de lengte van de as af. Bereken de extra
        afstand over de \qty{6000}{km} en de extra tijd die dat kost.
  \item De clou. Eén modern vezelpaar draagt ongeveer $10^{13}$ bits per
        seconde, en dit boek is ongeveer $2.4 \times 10^8$ bits. Bereken
        de heen-en-weertijd (``ping'') over de oceaan en het aantal
        exemplaren van dit boek dat de vezel per seconde verplaatst ---
        en zeg dan in één zin wat de wet van Snellius voor het internet
        doet.
\end{enumerate}
\end{problem}
